\begin{lemma}
\label{lemma-integral-over-invariants}
Let $S$ be a scheme. Let $(U, R, s, t, c)$ be a groupoid scheme over $S$.
Assume $U = \Spec(A)$ and $R = \Spec(B)$ are affine and
$s, t : R \to U$ finite locally free.
Let $C \subset A$ be as in (\ref{equation-invariants}).
Then $A$ is integral over $C$.
\end{lemma}

\begin{proof}
First, by Lemma \ref{lemma-finite-locally-free-disjoint-free}
we know that $(U, R, s, t, c)$ is a disjoint union
of groupoid schemes $(U_r, R_r, s, t, c)$ such that each $s, t : R_r \to U_r$
has constant rank $r$. As $U$ is quasi-compact, we have $U_r = \emptyset$ for
almost all $r$. It suffices to prove the lemma for each $(U_r, R_r, s, t, c)$
and hence we may assume that $s, t$ are finite locally free of rank $r$.

\medskip\noindent
Assume that $s, t$ are finite locally free of rank $r$.
Let $f \in A$. Consider the element $x - f \in A[x]$, where we think
of $x$ as the coordinate on $\mathbf{A}^1$.
Since
$$
(U \times \mathbf{A}^1, R \times \mathbf{A}^1,
s \times \text{id}_{\mathbf{A}^1},
t \times \text{id}_{\mathbf{A}^1},
c \times \text{id}_{\mathbf{A}^1})
$$
is also a groupoid scheme with finite source and target, we may apply
Lemma \ref{lemma-determinant-trick} to it and we see that
$P(x) = \text{Norm}_{s^\sharp}(t^\sharp(x - f))$
is an element of $C[x]$. Because $s^\sharp : A \to B$ is finite locally
free of rank $r$ we see that $P$ is monic of degree $r$.
Moreover $P(f) = 0$ by Cayley-Hamilton
(Algebra, Lemma \ref{algebra-lemma-charpoly}).
More precisely, Cayley-Hamilton says the polynomial
$s^\sharp(P) \in B[x]$ annihilates the element $t^\sharp(f)$.
Since the coefficients of $P$ are in $C$ we get that
$t^\sharp(P(f)) = 0$ in $B$. Since $t$ is faithfully flat,
we get $P(f) = 0$.
\end{proof}
